% Generated draft: requires mathematical and source review.
\documentclass[11pt]{article}
\usepackage[margin=25mm]{geometry}
\usepackage{fontspec}
\setmainfont{Latin Modern Roman}
\usepackage{amsmath,amssymb,mathtools,unicode-math}
\setmathfont{Latin Modern Math}
\usepackage{enumitem,needspace,xurl,hyperref}
\setlength{\parindent}{0pt}
\setlength{\parskip}{5pt}
\setlength{\emergencystretch}{3em}
\newcommand{\N}{\mathbb{N}}
\newcommand{\Z}{\mathbb{Z}}
\newcommand{\Q}{\mathbb{Q}}
\newcommand{\R}{\mathbb{R}}
\newcommand{\C}{\mathbb{C}}
\newcommand{\F}{\mathbb{F}}
\newcommand{\A}{\mathbb{A}}
\newcommand{\PP}{\mathbb{P}}
\newcommand{\E}{\mathbb{E}}
\newcommand{\eps}{\varepsilon}
\newcommand{\dd}{\,\mathrm{d}}
\begin{document}
\section*{UnsolvedMath ID: 30004722}
\textbf{Source problem code:} catalogue entry
\subsection{Definitions and mathematical statement}
For a smooth manifold \(M\), a Riemannian metric \(g\) has positive Ricci curvature if
\[
\operatorname{Ric}_g(v,v)>0
\quad\text{for every }x\in M\text{ and every }0\ne v\in T_xM.
\]
Write \(\mathcal R_{\operatorname{Ric}>0}(M)\) for the set of such metrics. Positive scalar curvature means \(\operatorname{Scal}_g(x)>0\) for every \(x\in M\).

To describe a conventional surgery comparison, let \(M\) be a closed connected smooth \(n\)-manifold and let
\[
e:S^k\times D^{n-k}\hookrightarrow M
\]
be a framed smooth embedding, where \(D^j\) denotes the closed \(j\)-disk. The manifold obtained by surgery on this embedded sphere is
\[
M_e=\left(M\setminus e\bigl(S^k\times\operatorname{int}D^{n-k}\bigr)\right)
\underset{S^k\times S^{n-k-1}}{\cup}
\left(D^{k+1}\times S^{n-k-1}\right),
\]
with gluing determined by the framing and with corners smoothed. The surgery sphere has codimension \(n-k\); the conventional codimension-at-least-three range is \(0\le k\le n-3\).

\textbf{Supplied research question.} Does a surgery theorem for positive Ricci curvature hold in the same generality as the Schoen--Yau and Gromov--Lawson surgery theorem for positive scalar curvature?

In the notation above, the existence-of-metrics conclusion being compared would take the form
\[
\mathcal R_{\operatorname{Ric}>0}(M)\ne\varnothing
\quad\Longrightarrow\quad
\mathcal R_{\operatorname{Ric}>0}(M_e)\ne\varnothing.
\]
\textbf{Unresolved scope.} The source does not explicitly quantify the dimensions, surgery indices, manifold class, or admissible embeddings, and does not state additional topological hypotheses or exclusions. The closed connected setting and codimension range above explain a conventional comparison; they are not confirmed hypotheses of the supplied question. In particular, quantifying over every \(n\ge3\), every \(0\le k\le n-3\), and every such framed embedding would be an additional interpretation, including \(k=0\), which is not asserted here as the intended open problem. The source also does not specify whether only existence of a metric is required or whether a geometric construction with prescribed behavior away from surgery is intended. Thus the phrase ``in the same generality'' remains an unresolved part of the formulation.
\subsection{Short English statement}
The question asks whether positive Ricci curvature admits a surgery theorem comparable in scope to the Schoen--Yau and Gromov--Lawson theorem for positive scalar curvature. The supplied record does not specify the precise allowed surgeries or extra hypotheses, so a fully quantified target requires clarification. Its open status is supplied dated curation, not an independently verified conclusion.
\subsection{Sources}
\begin{itemize}
\item Local UnsolvedMath registry, supplied complete source record: permanent ID \(30004722\), title \emph{Surgery Stability of Positive Ricci Curvature}, source problem code \texttt{catalogue entry}; collection: \emph{Oberwolfach Reports Open Problems}.
\item Supplied report citation: \emph{Analysis, Geometry and Topology of Positive Scalar Curvature Metrics} (2022), \emph{Oberwolfach Reports}, pp.\ 1579--1649, DOI \texttt{10.4171/owr/2021/30}. \url{https://doi.org/10.4171/owr/2021/30}. The record attributes the question to workshop open-problem discussions and supplies proposed year 2021, but names no individual proposer and gives no precise problem page within the report.
\item Supplied literature source: Reiser and Tripaldi, \emph{Calculus of Variations} (2026). \url{https://doi.org/10.1007/s00526-025-03211-2}.
\item Supplied literature source, labeled the open-access version of the surgery paper: \url{https://pmc.ncbi.nlm.nih.gov/articles/PMC12795942/}.
\item Literature assessment checked on \texttt{2026-08-22}: the supplied curation labels the problem open and reports that a positive Ricci curvature surgery theorem in the Schoen--Yau and Gromov--Lawson generality remains unknown, with only restricted surgery results available. This is supplied dated curation derived from the OWR record, not independently verified here; its time-sensitive status and cited claims require independent verification.
\end{itemize}
\end{document}
